%!TEX root = ../main.tex
% ANEXO I — Repositorios de código y material externo

\chapter{Repositorios de código fuente}\label{anx:repositorios}

\begin{guia}
Los anexos contienen material que no fue producido por el autor o que
es externo al documento: repositorios, datasets de terceros, normativa,
manuales del proveedor. Para cada repositorio: enlace, tecnología,
contenido y estado (activo, archivado).
\end{guia}

El código fuente de \sistema{} se aloja en un único repositorio público
de GitHub, organizado como monorepo.

\section{usal-tinku}

\prettygitlink{smercadocarbone/usal-tinku}

\textbf{Contenido.} El repositorio agrupa los tres componentes
desplegables del sistema, cada uno en su propia carpeta y con su propio
\texttt{README} de ejecución local:

\begin{itemize}
    \item \texttt{backend/}: \gls{api} REST en Java con Spring Boot,
          organizada como monolito modular.
    \item \texttt{matching-service/}: servicio de \gls{matching}
          semántico en Python, único proceso separado del monolito.
    \item \texttt{frontend/}: cliente \gls{pwa} en Next.js y React.
\end{itemize}

Se incluyen además la documentación de diseño (\texttt{docs/}), los
flujos de integración continua (\texttt{.github/workflows/}) y un
\texttt{docker-compose.yml} con la infraestructura local de desarrollo.

\textbf{Estado.} Activo, en desarrollo durante el cursado de la materia.

\completar{Licencia del repositorio (actualmente no declarada).}
